% NOTES:
% manually added german quotationmarks: \glqq \grqq // or use UTF-8 characters directly: „quotation marks“
% automatically added quotationmarks: \enquote{Text} // also supports nesting: \enquote{Test with \enquote{nested} quotation.}
% citation: \cite{Pohl2005} -> [1] /// \cite[S. 9ff]{Pohl2005} -> [2, S. 9ff]
%           (ISO-)Standards: http://itooktheredpill.dyndns.org/2010/citing-iso-standards-in-latex-using-bibtex/
% use of "et al.": http://blog.apastyle.org/apastyle/2011/02/et-al-when-and-how.html
%                  First occurence: Number of authors >= 6: use "et al.": "Sunderland et al. explained that.."
%                                   Number of authors <  6: name surnames of all authors
%                  Further occurence: Number of authors >= 3: use "et al.": "Sunderland et al. explained that.."
% references: Add label with \label{label1} and then refer to it with e.g. \hyperref[label1]{Label \ref*{label1} on page \pageref*{label1}}
%                                                                        or conveniently via \myRef{Abschnitt}{label}
% bibliography: keep information minimal (no URLs, no DOI, no ISBN)
% acronyms: \ac{KDE}  == K Desktop Environment (KDE) // on first occurence; on further use only acronym
%           \acs{KDE} == KDE // plural: \acsp{KDE} == KDEs
%           \acf{KDE} == K Desktop Environment (KDE) // \acfp{KDE} == K Desktop Environments (KDEs) 
%           \acl{KDE} == K Desktop Environment // plural: \aclp{KDE}
%       (!) \acp{KDE} == K Desktop Environments (KDEs) (on first occurrence) == KDEs (on further occurrences)
%           Definition: \acro{DD}{Data-Dictionary} // standard acronym definition
%                       \acrodefplural{DD}{Data-Dictionaries} // special plural definition (usually only an 's' is appended: \acp{KDE} == KDEs)
%                       Usage: \acl{DD} == Data-Dictionary // \aclp{DD} == Data-Dictionaries
% non-breaking space: used within abbreviations like z.\,B. to prevent both parts from being separated (e.g. due to line breaks)
%                     \, creates a narrow non-breaking space e.g. used for dates: 23.\,05.\,2042
%                     ~ creates a full non-breaking space e.g. used for full dates (23.\,Mai~2012)
%                        or terms that must not be separated (Prof.\,Dr.\,Bruno~Müller-Clostermann)
% wording: prefer the term "part" or "paragraph" to "chapter" unless you are writing a book or PhD thesis ;-) (German: "Abschnitt" statt "Kapitel" verwenden)
% LaTeX-Tabus: ftp://ftp.dante.de/tex-archive/info/german/l2tabu/l2tabu.pdf 
%
%
% PROBLEMS AND SOLUTIONS:
% Package hyperref Warning: Token not allowed in a PDFDocEncoded string... removing `math shift' on input line XXX.
%    Solution: provide alternative plain text for mathsymbols or special characters that will end up in links. 
%    E.g. change \section{The joy of $\pi$} to \section{The joy of \texorpdfstring{$\pi$}{pi}}
%    http://homepages.inf.ed.ac.uk/imurray2/compnotes/latex.html
% Biliography is not created properly
%    Problem: References are not found after building, log file and \printbibliography command show messages such as "empty bibliography".
%    Try to run biber manually on command line. Possible biber caching error (encountered on OS X 10.8.2):
%    "data source /var/folders/m6/bn7r45zx6cx55rr9g4qh6s6w0000gn/T/par-agoldst/<..>/inc/lib/Biber/LaTeX/recode_data.xml not found in ."
%    Solution: sudo rm -rf /var/folders/m6/bn7r45zx6cx55rr9g4qh6s6w0000gn/T/par-agoldst/
%    see http://humtex.wordpress.com/2011/11/29/biber-first-aid-for-data-source-not-found/
%
%

% ########################### DOCUMENTCLASS ###########################

\documentclass[paper=a4, % equivalent to a4paper
               headsepline, % line after headline
               fontsize=12pt,
               headings=small,
               %draft, % debug mode
               oneside, % one sided document: no alternating page number positions and margins (equivalent to twoside=false)
               ngerman, % new German spelling rules (important for hyphenation)
               captions=tableheading, % within table environment: put \caption{} text above table
               toc=bibliographynumbered, % adds bibliography to table of contents (numbered like chapters; unnumbered with toc=bibliography)
               1.5headlines]
               {scrreprt}

% ######################### PACKAGES + CONFIG #########################

\usepackage[T1]{fontenc} % German umlauts, special chars, etc.
\usepackage[utf8]{inputenc} % input encoding (preferably utf8 with utf8 encoding within your editor! ansinew, latin1, etc. are also possible but require you to use cumbersome workarounds to write German umlauts e.g. "u for writing an ü)
\usepackage[ngerman]{babel} % new German spelling rules (important for hyphenation)
\usepackage[german]{varioref} % for commands like \vref to create references like "...see figures 3.4 to 3.6 on pages 23–24..."


% special symbols
\usepackage{pifont}

% strings have to be entered after utf8 has been set as input encoding

% strings for the frontpage
\newcommand{\documentType}{Bachelorarbeit} % Art der Arbeit: Master-/Bachelor-/Seminararbeit/etc.
\newcommand{\documentHeading}{Natural Language Processing zur Verbesserung von Suchfunktionen in der Software-Entwicklung}
\newcommand{\documentLocation}{Essen} % also used as location for declaration of academic honesty / Ort in Eidesstattlicher Erklärung
\newcommand{\documentDate}{\today}%10.\,04.\,2014} % also used as date for declaration of academic honesty / Datum in Eidesstattlicher Erklärung

% PDF document metadata (the following strings must not contain line breaks, non-breaking spaces, etc.)
\newcommand{\documentPDFauthor}{Leon Zimmermann}
\newcommand{\documentPDFtitle}{\documentHeading}



\usepackage{acronym} % enables acronym environment ftp://ftp.tu-chemnitz.de/pub/tex/macros/latex/contrib/acronym/acronym.pdf
                     % http://de.wikibooks.org/wiki/LaTeX-W%C3%B6rterbuch:_Abk%C3%BCrzungsverzeichnis

\usepackage{hyphenat} % hyphenation
\hyphenation{} % special hyphenation for certain words (you can also simply add words that must not be hyphenated)

\usepackage{eurosym} % Euro currency symbol: use \EUR{10} or 10 \euro (\EUR{10} recommended because of correct space between number and symbol)
\DeclareUnicodeCharacter{20AC}{\euro} % UTF-8: enables use of Euro currency sign € (cf. http://www.theiling.de/eurosym.html)

\usepackage{amsmath, amsthm, amssymb} % AMS math symbols and theorem boxes; configure via \newtheoremstyle (cf. ftp://ftp.ams.org/pub/tex/doc/amscls/amsthdoc.pdf)
\usepackage{mathtools}

\usepackage[top=2.5cm, bottom=2cm, left=5cm, right=1.4cm, headsep=1cm, footskip=1cm]{geometry} % borders
%\renewcommand{\restoregeometry}{\newgeometry{top=2.5cm, bottom=2cm, left=5cm, right=1.4cm, headsep=1cm, footskip=1cm}}
    % examples from http://www.artofproblemsolving.com/Wiki/index.php/LaTeX:Layout
    % headsep  = Distance from bottom of header to the body of text on a page.
    % footskip = Distance from bottom of body to the bottom of the footer.
\usepackage{scrlayer-scrpage} % header/footer https://ctan.org/pkg/scrlayer-scrpage?lang=de

\usepackage{setspace} % enables line spacing via \singlespacing, \onehalfspacing, \doublespacing or \setstretch{<factor>} for custom spacing
\usepackage{verbatimbox}
\usepackage{rotating} % rotate objects e.g. with \rotatebox{90}{}; also used for tables in landscape format with sidewaystable environment
\usepackage{tabularx}
\usepackage{calc}
\usepackage{array}
\usepackage[font={bf, small, sf}]{caption} % customize the captions in floating environments
\usepackage{subcaption}
\usepackage[final]{pdfpages} % inclusion of external PDF documents (options e.g.: selected pages including a blank one: pages={3,{},8-11,15} )
\usepackage{multirow} % span column over multiple rows
\usepackage{booktabs} % table rules with various thicknesses (\toprule, \midrule, \bottomrule, etc.)
\usepackage{longtable, tabu} % tables that continue to the next page; http://mirror.unl.edu/ctan/macros/latex/contrib/tabu/tabu.pdf 
\usepackage{ltcaption} % caption for longtable
\usepackage{multicol} % add areas with multiple columns
\usepackage{enumitem} % prefix for enumerations
\setlistdepth{9} % enumitem: enable long nested lists
\setlist[itemize,1]{label=$\bullet$}
\setlist[itemize,2]{label=$\bullet$}
\setlist[itemize,3]{label=$\bullet$}
\setlist[itemize,4]{label=$\bullet$}
\setlist[itemize,5]{label=$\bullet$}
\setlist[itemize,6]{label=$\bullet$}
\setlist[itemize,7]{label=$\bullet$}
\setlist[itemize,8]{label=$\bullet$}
\setlist[itemize,9]{label=$\bullet$}
\renewlist{itemize}{itemize}{9}

%\bibliographystyle{alpha} % examples: http://amath.colorado.edu/documentation/LaTeX/reference/faq/bibstyles.html
%\bibliographystyle{dinat} % bibliographies following the german norm DIN 1505
\usepackage[bibstyle=numeric,
            citestyle=numeric,
            backend=biber,
            maxcitenames=2,
            maxbibnames=100,
            sortcites=true,
            doi=false,
            isbn=false,
            url=false]{biblatex} % bibliography package; replaces old natbib package; numeric = Reference in brackets [1].
                                        % remember to translate via "biber file.bcf" (or configure LaTeX editor)
                                        % very useful for printing parts of the bibliography within the document
                                        % see http://www.komascript.de/node/1220 for filter examples
                                        % http://tex.stackexchange.com/questions/32851/biblatex-combine-filters
                                        % http://biblatex.dominik-wassenhoven.de/download/DTK-2_2008-biblatex-Teil1.pdf
%\nocite{*} % display all entries in bibliography regardless of whether they have been cited or not
\DefineBibliographyStrings{ngerman}{andothers={et~al\adddot}} % "u.a." zu "et al."

\usepackage[german=quotes]{csquotes} % dynamic and configurable quotationmarks that \enquote{also support \enquote{nested} quotations.}


\usepackage{lmodern} % Latin Modern fonts
%\usepackage{picins} % insert pictures into paragraphs

\usepackage{float}
\usepackage{wrapfig} % wrap text around figures with wrapfigure environment (has to be loaded right after float package)
\restylefloat{figure} % float only at the place where you define it (without LaTeX moving it)
\floatstyle{plain} % standard: no box is drawn


\usepackage{color} % define own colors: \definecolor{MyColor}{rgb}{0.37,0.17,0.97}
\usepackage{xcolor}
\usepackage{colortbl} % color tables and cells
\usepackage{adjustbox} % http://mirror.unl.edu/ctan/macros/latex/contrib/adjustbox/adjustbox.pdf

\usepackage{framed} % framed or shaded regions that can break across pages

\usepackage{graphicx} % enhanced support for graphics

\usepackage{url} % e.g. \url{http://www.example.com}

\usepackage{fullwidth} % full width environment to display something beyond the text width http://www.ctan.org/tex-archive/macros/latex/contrib/fullwidth


% hyperref package has to be loaded after the float package
\usepackage[%pdftex=true, % only if you use pdftex
            hypertexnames=true,
            pdfstartview={FitH}, % fit the width of the page to the window
            colorlinks=true,
            linkcolor=black,
            citecolor=black,
            filecolor=black,
            urlcolor=black]{hyperref} % clickable urls in PDF: \href{URL}{text}
\hypersetup{% configure PDF metadata with \hypersetup command to pass UTF-8 characters (e.g. umlauts) properly
            pdftitle={\documentPDFtitle}, % title (PDF metadata)
            pdfauthor={\documentPDFauthor}, % author (PDF metadata)
}
\urlstyle{tt} % font style of url: rm = CM Roman; sf = CM Sans Serif; tt = CM Typewriter; same = selected font


% Pseudo Sourcecode Formatting with algorithmicx (http://ctan.org/pkg/algorithmicx)
%   Packages have to be loaded after float and hyperref
%   Alternative: algorithm2e which cannot split code on multiple pages but has nice vertical lines
%   Get vertical lines in algorithmicx as seen in algorithm2e: http://tex.stackexchange.com/questions/110431/ploblems-with-vertical-lines-in-algorithmicx
%   Chaining procedure calls: http://tex.stackexchange.com/questions/16046/how-to-nest-call-in-algorithmicx
\usepackage{algpseudocode} % part of algorithmicx
\usepackage[Algorithmus]{algorithm} % float environment algorithm, parameter is the prefix title in every algorithm box (e.g. "Algorithm 1")


% margin notes
\setlength{\marginparwidth}{1.2in} % width of margin notes
\reversemarginpar % margin notes (added via \marginpar{margin text}) are placed on the opposite side (left/inside instead of right/outside)


% load custom commands and strings
% suppresses page numbering in the header for one page (usage e.g.: \suppressPagenumber{\part{Stand der Forschung}})
\newcommand{\suppressPagenumber}[1]{{\ohead[]{} #1 \ohead[\rightline{\thepage}]{\rightline{\thepage}}}}


% own quotation style for a big single quotation within a block of text
\newcommand{\myQuotation}[2]{\begin{quote}\small{\textit{\enquote{#1}}} -- #2\end{quote}}



% side notes with marginpar
\newcommand{\myNote}[1]{\marginpar{\raggedleft \footnotesize \color{gray} \textit{#1}}}



% cite reference with full title ahead
\newcommand{\myCiteTitle}[1]{\citetitle{#1} (\citeauthor{#1}~\cite{#1})}




% Image/Figure
% Usage: \myImg{caption}{label}{width-factor}{path}
%   Example: \myImg{Beispiel eines Datenflussdiagramms}{fig:dfd}{1}{data-flow-diagram.pdf}
% http://www2.informatik.hu-berlin.de/~piefel/LaTeX-PS/Archive-2004/V08-grafikein.pdf
\newcommand{\myImg}[4]{
    \begin{figure}[h!]
        \centering
        \includegraphics[width=\textwidth * \real{#3}]{./img/#4}
        \caption{\label{#2}#1}
    \end{figure}    
}
\newcommand{\myImgCfg}[5]{
    \begin{figure}[#5]
        \centering
        \includegraphics[width=\textwidth * \real{#3}]{./img/#4}
        \caption{\label{#2}#1}
    \end{figure}    
}

% Shortcut for "..Abbildung 2 zeigt.." via "..\myImgRef{fig:screenshot} zeigt.."
\newcommand{\myImgRef}[1]{\hyperref[#1]{Abbildung~\ref*{#1}}}



% Shortcut for "..Abschnitt 2 beschreibt.." via "..\myRef{Abschnitt}{chap:myrefchapter}.."
\newcommand{\myRef}[2]{\hyperref[#2]{#1~\ref*{#2}}}


% Shortcut for "..auf Seite 7.." via "..\myPageRef{fig:myfigure}.."
\newcommand{\myPageRef}[1]{\hyperref[#1]{Seite~\pageref*{#1}}}



% Shortcuts for frequently used strings
\newcommand{\sczb}{z.\,B. } % zum Beispiel; added space because those phrases will never be used at the end of a sentence
\newcommand{\scdh}{d.\,h. } % das heißt
\newcommand{\scua}{u.\,a. } % unter anderem






% Mathematics
\newcommand{\powerset}[1]{\mathcal P~\left({#1}\right)}
%\newcommand{\powerset}[1]{\wp \left({#1}\right)}



% Font
% changes font family ptm = Times, phv = Helvetica,.. http://de.wikibooks.org/wiki/LaTeX-W%C3%B6rterbuch:_fontfamily
%\renewcommand{\familydefault}{phv}



% colored box
% see LaTeX - fcolorbox.png
% http://www.lessjunkmorefunk.de/de/node/10
%\definecolor{Gray}{gray}{0.9}
%\newcommand{\inhalt}[1]{\fcolorbox{black}{Gray}{\parbox{\textwidth}{#1}}}


\addbibresource{./referenzen.bib} % ./ for local folder is very important!


% ############################# DOCUMENT ##############################
\begin{document}
\shorthandoff{"} % deactivates shorthand "a for a umlaut (ä) (necessary if quotation marks are used e.g. in source code)
\setcounter{secnumdepth}{3} % numbering up to three level
%\sloppy % LaTeX will always try to create justified text (this may result in bigger gaps between words)

% ############################# FRONTPAGE #############################
\pagestyle{empty} % remove page number

\includegraphics[width=\textwidth]{./img/logo_unidue_paluno_deutsch.pdf}

\begin{sffamily}

\begin{center}

\vspace*{0.7cm}

\normalsize{Institut für Informatik und Wirtschaftsinformatik (ICB)\\%Institute for Computer Science and Business Information Systems (ICB)
Lehrstuhl für Software Engineering, insb. mobile Anwendungen\\
Prof. Dr. Volker Gruhn}

%\vspace{\baselineskip}

\vspace*{2.5cm}

\LARGE{\documentHeading}

\vspace*{1cm}

\large{\documentType}

\vspace*{2.5cm}

\normalsize{vorgelegt der Fakultät für Wirtschaftswissenschaften\\
der Universität Duisburg-Essen (Campus Essen) von}

\vspace*{\baselineskip}

\normalsize{
Leon Zimmermann \\
Laddringsweg 8 \\
45219 Essen \\
Matrikelnummer: 3080384
}

\vspace*{1.5cm}

\small{\documentLocation, den \documentDate}

\vspace*{1.5cm}

\normalsize{
\begin{tabular*}{\linewidth}{p{0.3\linewidth}p{0.7\linewidth}}
\textbf{Betreuung:}      & Wilhelm Koop, Sascha Feldmann\\
\textbf{Erstgutachter:}  & Prof.\,Dr.\,Volker~Gruhn\\
\textbf{Zweitgutachter:} & ???\\
\\
\textbf{Studiengang:}    & Angewandte~Informatik - Systems~Engineering (B.\,Sc.)\\
\textbf{Semester:}       & 10
\end{tabular*}
}

\end{center}

\end{sffamily}
\clearpage % calls \pagestyle{empty}
%\onehalfspacing % set one-and-a-half line spacing after frontpage was included
\setstretch{1.1} % a value of 1.25 refers to one-and-a-half spacing and 1.667 to double-spacing (for 10pt)
                 % (LaTeX calculates with a leading of 1.2 at a point size of 10pt: 1.2 * 1.25 = 1.5 line spacing)

% ############################# HEADLINES ######################
\renewcommand{\headfont}{\small} % small headline font size to fit even long headlines onto the page

\clearscrheadfoot
\ohead[\rightline{\thepage}]{\rightline{\thepage}} % headline with number of page on the right hand side
\automark{chapter}

\pagenumbering{roman} % roman numbering: I, II, III, IV, ...

\pagestyle{scrheadings} % defines whether headlines are printed and how they are styled

% ############################## ABSTRACT #############################

%- Was ist das generelle Thema der Arbeit?
%Beispiel: Kontext ist der Schnittpunkt Variabilitätsmodellierung / SPLE / SA.
%
%- Was ist das ganz konkrete Problem, das in der Arbeit behandelt wird? Was ist die Forschungslücke?
%Beispiel: DFDs haben keine Variabilitätserweiterung. Bisherige Lösungen behandeln das Thema nur abstrakt und können eine Konsistenzprüfung im Sinne des SPLE-Gedankens nicht konkret umsetzen.
%
%- Was ist die generelle Vorgehensweise (Methodik, z.B. Art der Studie: Experiment, Literatursuche, …)?
%Beispiel: SLR zum Aufspüren von Problemen und potentiellen Lösungen mit dem Ziel, diese zu einem eigenen Ansatz zu kombinieren.
%
%- Was ist das zentrale Resultat? Was folgt aus dem Resultat (Wert der Arbeit, z.B. für künftige Studien)?
%Beispiel: Ansatz zur SA im SPLE mit Toolunterstützung. Macht den Weg frei für weiterführende Ansätze in dem Rahmen (evtl. ein knackiges Beispiel nennen, was man tun kann).
%
%Dabei sollte die Wortwahl – im Gegensatz zu den von mir gegebenen Stichpunkten – natürlich durchgehend knackig und äußerst präzise sein. Zu jedem Punkt sollten jeweils ein bis zwei Sätze formuliert werden.
%Insgesamt sollte ein Abstract ca. 100-250 Wörter lang sein (grob eine eine halbe Seite bei dem üblichen Format, niemals mehr als eine Seite).
%
%Du solltest dabei daran denken, dass ein Abstract potentiellen Lesern zur Grobeinordnung Deiner Arbeit dient, z.B., um festzustellen, ob die Arbeit für ihre Forschung relevant sein könnte. Außerdem ist ein Abstract das Aushängeschild der Arbeit; er sollte bei am Problem interessierten Lesern (auch bei Gutachtern!) großes Interesse wecken.
%%
%
%
%
\section*{Abstract}
% TODO: Zusammenfassung auf Englisch

%
%
%
\section*{Zusammenfassung}
Softwareentwickler müssen sich bei ihrer Arbeit Informationen zusammensuchen, welche sie für die weitere Arbeit benötigen.
So muss ein Softwareentwickler bei der Implementierung eines Features die Intention des Features kennen.
Solche Informationen können in Wissensdatenbanken hinterlegt sein.
Um dort die gewünschten Informationen zu finden können Suchfunktionen verwendet werden.
Insbesondere, wenn dem Softwareentwickler im Vorfeld nicht klar ist, wo er die gewünschten Informationen in der Wissensdatenbank finden kann.
Aber nicht immer liefert diese Suchfunktion die gewünschten Informationen.\\

In dieser Arbeit wird eine semantische Suchfunktion entwickelt.
Diese soll eine Verbesserung zu bestehenden Suchfunktionen bieten, wie beispielsweise die Suchfunktion von Confluence.
Um festzustellen, ob die neue Suchfunktion \textit{besser} ist als die bestehende, werden Methoden und Kriterien zur Evaluierung von Suchfunktionen erörtert.
Anhand dieser Kriterien werden die neu implementierte Suchfunktion und die bestehende Suchfunktion von Confluence in einer Studie verglichen.
In der Studie werden beispielhafte Sucheingaben definiert, sowie die erwarteten Ergebnisse.
Die Definition der Sucheingaben erfolgt auf Basis von Anwendungsfällen.
Die Anwendungsfälle beschreiben Situationen, in denen es realistisch ist, dass ein Softwareentwickler die Suchfunktion einer Wissensdatenbank verwendet.
Anhand von Argumentationen werden entsprechend realistische Sucheingaben erstellt.\\

% TODO: Ergebnis der Studie darstellen

% ########################### TOC + ACRONYMS ##########################
\setcounter{chapter}{0}

\setcounter{tocdepth}{2} % table of contents depth: 3 means e.g. 4.3.2.1 will be listed, whereas 2 will lead to 4.3.2. as the deepest entry, etc.
{
\sloppy
\tableofcontents 
}

\listoffigures

\listoftables

\chapter*{Abkürzungsverzeichnis}

\begin{acronym}
    \acro{ANN}{Approximate Nearest Neighbor}
    \acro{HNSW}{Hierarchical Navigable Small Worlds}
    \acro{NLP}{Natural Language Processing}
    \acro{VSM}{Vector Space Model}
\end{acronym}
\clearpage % Clear page to get the arabic page counting on the next page

% ########################## LAYOUT + DESIGN ##########################

\clearscrheadfoot % cleans footer from previous configurations
\ihead[]{\leftline{\leftmark}} % current section and chapter heading in header on the left
\ohead[\rightline{\thepage}]{\rightline{\thepage}} % page number in header on the right
\automark{chapter} % for creating the current section/chapter in the heading

\pagenumbering{arabic} % arabic numbering: 1, 2, 3, 4, ...
\setcounter{page}{1}

\addtokomafont{chapter}{\LARGE}
\addtokomafont{section}{\Large}
\addtokomafont{subsection}{\large}
\addtokomafont{paragraph}{\itshape}

\renewcommand{\thepart}{}%\Alph{part}}
\renewcommand{\thesubsubsection}{\Roman{subsubsection}}

\renewcaptionname{ngerman}{\partname}{}

% ############################## CONTENT ##############################

\chapter{Einleitung}
Für die tägliche Arbeit benötigt ein Softwareentwickler Informationen, welche über den Code, an dem er arbeitet, hinausgehen.
Um an diese Informationen zu kommen kann der Softwareentwickler eine Person suchen, welche ihm die gewünschte Information geben kann.
Diese Person kann er im Projekt finden, oder über Websites, wie StackOverflow.
Darüber hinaus wird in vielen Projekten eine Wissensdatenbank angelegt, welche Informationen enthält, welche spezifisch für das Projekt sind.
Eine solche Wissensdatenbank ist Confluence.
Sie bietet eine Suchfunktion, welche es dem Softwareentwickler erleichtern soll, die gewünschten Informationen zu finden.
Beispiele für Informationen sind die Spezifikationen oder Dokumentationen eines Teiles der Software, mit welcher der Softwareentwickler gerade arbeitet.
Oder aber auch Best-Practices, Guides, oder Informationen darüber, wie die Software gestartet oder ausgeliefert wird.
Auch Informationen über den Projektplan sind für einen Softwareentwickler von Bedeutung.\\

Aber nicht immer finden Softwareentwickler die gewünschten Informationen mithilfe der Suchfunktion.
Ziel dieser Arbeit soll es sein, mithilfe von Wissen über Suchalgorithmen und Algorithmen aus dem Natural Language Processing zu zeigen, wie sich bestehende Suchfunktionen verbessern lassen.
Software, wie chatGPT\footnote{https://openai.com/blog/chatgpt} zeigt, dass Large Language Models eine valide Möglichkeit zur Information Extraction sind.
Es ist denkbar, dass Suchfunktionen für Softwareentwickler durch Large Language Models verbessert werden können.
Diese Arbeit erörtert, wie die Technologien von Large Language Models verwendet werden können, um eine semantsiche Suche zu implementieren.
Die implementierte semantische Suche wird mit der bestehenden Confluence-Suche verglichen, um herauszufinden, ob die semantische Suche eine Verbesserung zu der bestehenden Confluence-Suche darstellt.

\section{Vorgehensweise}
Die Gründe, warum eine gewünschte Information schwierig zu finden ist, sind vielzählig.
Manchmal kennt der Softwareentwickler nicht das genaue Wording, um die gewünschten Informationen zu finden.
Manchmal ist das Abstraktionslevel der gefundenen Informationen nicht das, welches sich der Softwareentwickler gewünscht hat.
Beispielsweise, wenn eine allgemeine Definition von Domänenobjekten gesucht wird, aber eine Spezifikation eines Anwendungsfalls gefunden wird, in welchem das Domänenobjekt lediglich erwähnt wird.
Im Rahmen dieser Arbeit wird eine neue Suchfunktion entwickelt, welche eine Verbesserung zu bestehenden Suchfunktionen von Wissensdatenbanken, wie Confluence, darstellen soll.
Nachdem die neue Suchfunktion entwickelt wurde, wird diese mit einer bestehenden Suchfunktion verglichen, um zu untersuchen, ob die neue Suchfunktion besser ist als eine bestehende Suchfunktion.
Dazu werden objektive Metriken zur Messung von Suchfunktionen dargestellt.
Die einzelnen Schritte für die Entwicklung und Evaluierung der neuen Suchfunktion werden im Folgenden dargestellt:
\begin{itemize}
   \item \textbf{Definition von Anwendungsfällen:}
         Es werden zuerst Anwendungsfälle definiert.
         Die Anwendungsfälle beschreiben die konkreten Situationen, in welchen ein Softwareentwickler eine Suche nutzen könnte.
         Das hilft später dabei Suchfunktionen miteinander zu vergleichen.
         Denn anhand der Anwendungsfälle können realistische Sucheingaben definiert werden.
         Diese Sucheingaben können an zwei verschiedene Suchfunktionen übergeben werden.
         Anschließend können die gefundenen Ergebnisse verglichen werden.
         Die genaue Vorgehensweise für diesen Vergleich wird in \myRef{Kapitel}{chap:vergleich-der-suchfunktionen} erklärt.
   \item \textbf{Technologien von Confluence:}
         Als nächstes wird die Technologie von Confluence dargestellt.
         Das hilft dabei später die Konzeption einer neuen Suchfunktion zu entwickeln und zu begründen, warum die neue Technologie besser sein soll.
         Im weiteren Verlauf der Arbeit wird erläutert, dass für den Vergleich der Suchfunktionen nicht die tatsächliche Confluence-Suchfunktion als Benchmark für die neue Suchfunktion herangezogen wird.
         Stattdessen wird die neue Suchfunktion so entwickelt, dass sie so konfiguriert werden kann, dass sie eine gute Heuristik für die Confluence-Suchfunktion darstellt.
         Dieses Kapitel erfüllt ebenfalls den Zweck eine vergleichbare Suchfunktion entwickeln zu können, welche anhand der objektiven Kriterien gemessen werden kann.
         Auch bei der Auswertung der Ergebnisse der Suchfunktion sind die Inhalte aus diesem Kapitel relevant.
   \item \textbf{Konzeption der neuen Suchfunktion:}
         Nachdem die Technologien von Confluence erörtert wurden, wird in diesem Kapitel die Theorie für die Implementierung einer Suchfunktion erläutert.
         Es wird zuerst erklärt, dass die neue Suchfunktion eine semantische Suchfunktion ist.
         Dann wird erläutert, was eine semantische Suche ist, und warum eine semantische Suche eine Verbesserung gegenüber den Technologien von Confluence darstellen soll.
         Anschließend wird erklärt, welche Technologien für die Implementierung einer semantischen Suche notwendig sind.
         Zuletzt wird die Implementierung der semantischen Suche mithilfe von der Vektordatenbank Weaviate\footnote{https://weaviate.io/} erklärt.
   \item \textbf{Vergleich der Suchfunktionen:}
         In diesem Kapitel werden zuerst Methoden für die Evaluierung von Suchfunktionen herausgearbeitet.
         Damit wird die Frage beantwortet, wann eine Suchfunktion \textit{gut} ist.
         Das hier erläuterte Wissen wird für die Durchführung der Studie benötigt.
         Denn um festzustellen, ob die Implementierung eine Verbesserung darstellt, muss eine Studie durchgeführt werden.
         Aufgrund des Scopes der Arbeit wird nur eine rudimentäre Studie durchgeführt.
         Die Studie vergleicht die Precision-Werte der Confluence-Suchfunktion und der neuen Suchfunktion mithilfe eines generierten Datensatzes von Sucheingaben und erwarteten Ergebnissen.
         Grundlage für die Generierung dieses Datensatzes sind die Anwendungsfälle, welche zuvor definiert wurden.
         Die Ergebnisse werden dargestellt und diskutiert.
         Es wird gezeigt, dass eine semantische Suche für die generierten Daten und den spezifischen Kontext der Datengrundlage, einen schlechten Precision-Wert besitzt. 
         Dann werden dieses Ergebnis interpretiert und es wird der Schluss gezogen, dass Sentence Transformer besser geeignet sind für Open-Domain Datensätze, im Gegensatz zu Closed-Domain Datensätzen.
         Zuletzt wird der Studienaufbau diskutiert.
    \item \textbf{Zusammenfassung und Ausblick}
         In diesem Kapitel werden die einzelnen Kapitel der Arbeit nochmal zusammenfassend dargestellt.
         Außerdem wird das Ergebnis der Arbeit zusammenfassend dargestellt.
         In \myRef{Kapitel}{chap:ausblick} werden Ansätze erörtert, welche die Suchergebnisse von semantischen Suchen in spezifischen Domänen verbessern sollen.
\end{itemize}

\section{Verwandte Arbeiten}
Es gibt einige Arbeiten, welche die gleichen oder sehr ähnliche Probleme adressieren.
Die verschiedenen Herangehensweisen werden in diesem Kapitel erläutert.
Zuerst wird der Begriff Traceability behandelt, und wie diese genutzt wird, um Suchfunktionen zu verbessern.
Die Herangehensweisen zur Herstellung von Traceability umfassen die Verwendung von Word Embeddings oder auch ein Task-Based Ansatz.
Es folgt eine Übersicht über Arbeiten, welche sich mit Question Answering beschäftigen.
Dann folgt eine Übersicht über Arbeiten, welche sich mit der Evaluierung von Suchfunktionen auseinandersetzen.

% Traceability
Traceability bedeutet, dass sich von einer Stelle im Code, auf die entsprechenden Stellen in anderen Artefakten zurückschließen lässt.
Ein Anwendungsfall für eine solche Traceability-Funktionalität ist \textit{Feature Location}, also das Finden der Spezifikation eines Features, wenn nur der Code vorhanden ist.
Analog dazu ist \textit{Bug Localization} ein Anwendungsfall zum Auffinden von Stellen im Code, welche mit einem Bug zusammenhängen.
Zur Herstellung von Traceability zwischen Code und Dokumentation gibt es verschiedene Ansätze.
Haiduc et. al. schlagen ein System zur Verbesserung von Sucheingaben vorgeschlagen (\cite{Haiduc_Bavota_Marcus_Oliveto_DeLucia_Menzies_2013}).
Das System verwendet Query Reformulations, um die Traceability herzustellen.
Query Reformulation bedeutet, dass das System den Softwareentwickler bei der Eingabe einer Suchanfrage zur Suche nach den passenden Artefakten unterstützt.
Dazu gibt der Softwareentwickler zunächst eine Suchanfrage ein, und markiert diejenigen Ergebnisse, welche am relevantesten für ihn sind.
Auf Grundlage der gewählten Ergebnisse und mithilfe eines Machine Learning Algorithmus werden nun Vorschläge für eine verbesserte Suchanfrage gemacht.
Dabei gibt es verschiedene Strategien.
Wenn der Softwareentwickler zu Beginn eine sehr lange Suchanfrage eingegeben hat, dann kann das System eine Reduktion der Suchanfrage vorschlagen.
Hat der Softwareentwickler dagegen lediglich einen Suchbegriff angegeben, so kann das System eine Erweiterung der Suchbegriffe vorschlagen.
Dazu greift das System auf Synonyme des eingegebenen Suchbegriffes zurück.\\

Ye et. al. beschreiben, wie Word Embeddings dazu verwendet werden können, um Traceability zwischen Code und anderen Softwareentwicklungs-Artefakten herzustellen.
Word Embeddings sind eine Datenstruktur, welche einem Wort einen Vektor in einem n-dimensionalen Raum zuweist.
Anhand dieses Vektors kann die Ähnlichkeit zwischen Wörtern beschrieben werden.
Ähnliche Wörter haben eine geringe Distanz im n-dimensionalen Raum.
Unähnliche Wörter haben eine hohe Distanz.
Der Algorithmus, welcher die Ähnlichkeit der Wörter bestimmt, macht Gebrauch von der Distributional Hypothesis (\cite{Harris_1954}).
Dieser besagt, dass Wörter, welche im gleichen Kontext verwendet werden, eine ähnliche Semantik besitzen.
Hermit wird also die Ähnlichkeit der Wörter bestimmt.
Dieses Verfahren wird nun sowohl auf den Code angewendet als auch auf die Softwareentwicklungs-Artefakte (\cite{Ye_Shen_Ma_Bunescu_Liu_2016}).\\

Antoniol et. al. verwenden einen ähnlichen Ansatz, wie Ye et. al.
Auch hier werden Word Embeddings verwendet um Softwareentwicklungs-Artefakte gegen den Code zu matchen.
Hier durchlaufen die Artefakte und der Code zwei verschiedene Pipelines.
Die Wörter der Artefakte in natürlicher Sprache werden in lowercase umgewandelt.
Anschließend werden Stoppwörter entfernt.
Zuletzt werden Flexionen entfernt.
Aus dem Code werden zunächst Identifier extrahiert.
Identifier, welche mehrere Wörter unter Verwendung von CamelCase oder snake\_case beinhalten, werden in die einzelnen Wörter aufgeteilt.
Anschließend werden die Identifier auf die gleiche Art und Weise normalisiert, wie die Wörter der Softwareentwicklungs-Artefakte.
Dann erfolgt sowohl für die Identifier als auch für die Wörter aus den Artefakten die Indizierung, also die Umwandlung in Word Embeddings (\cite{Antoniol_Canfora_Casazza_DeLucia_2000}).\\

Treude et. al. entwickeln eine Oberfläche, welche die Suche von \textit{Tasks} ermöglicht.
Damit soll ebenfalls Traceability zwischen Code und Dokumentation hergestellt werden.
Dabei ist unter Task eine Operation im Code zu verstehen.
Sie beschreiben einen Task als Verben, welche mit einem direkten Objekt oder einer Präposition in Verbindung stehen.
Die Autoren nennen die Phrasen \textit{get iterator} und \textit{get iterator for collection} als Beispiele.
Die Software analysiert nun die gesamte Dokumentation und extrahiert Tasks.
Die Tasks werden in einen Index geschrieben, sodass der Softwareentwickler nach ihnen suchen kann (\cite{Treude_Sicard_Klocke_Robillard_2015}).\\

% Question Answering
Neben der Traceability zwischen Code und Dokumentation ist das Question Answering ein Ansatz zur Verbesserung von Suchfunktionen.
Suchfunktionen sind Document Retrieval Systeme.
Sie liefern Dokumente, welche zu der Sucheingabe des Nutzers passen.
Document Retrieval Systeme sind eine Unterkategorie von Information Retrieval Systemen.
Information Retrieval Systeme liefern auf Anfrage Informationen an den Nutzer.
Im Fall einer Suchfunktion werden Dokumente geliefert, welche diese Information beinhalten.
Mithilfe der Dokumente ist der Ort, an dem sich die, vom Nutzer gewünschte, Information befindet eingegrenzt.
Nichtsdestotrotz muss der Nutzer aus dieser Eingrenzung die gewünschte Information manuell extrahieren.
Ganz im Gegensatz zu Question Answering Systemen.
Zhang et. al. untersuchen verschiedene Herangehensweisen zur Implementierung von Open-Domain Question Answering Systemen (\cite{Zhang_Chen_Xu_Cao_Chen_Cohn_Fang_2023}).
Open-Domain Question Answering Systeme beantworten allgemeine Fragen eines Nutzers, z.B. basierend auf Informationen von Wikipedia.
Closed-Domain Question Answering Systeme beantworten dagegen Fragen im Kontext einer spezifischen Domäne, z.B. basierend auf unternehmensinternen Informationen.\\
% TODO: closed-domain sentence transformer transformer zu verwandten arbeiten hinzufügen (fine-tuned transformers)

% TODO: Informationen lassen sich auch anreichern: NER etc. (NLP Techniken)
% TODO: Informationen lassen sich aufbereiten: RAG

% Messung von Performance
% TODO: Trec und Benchmarks hinzufügen
Für die Messung der Performance einer Suchfunktion gibt es verschiedene Metriken.
Metriken können systemspezifisch sein, nutzerspezifisch.
Systemspezifische Metriken messen die Performance anhand objektiver Kriterien, welche an dem System gemessen werden können.
nutzerspezifische Metriken messen die Performance dagegen anhand subjektiver Kriterien.
Die Messung erfordert Probanden, welche die Suchfunktion verwenden.
Die Probanden legen die Performance der Suchfunktion anhand der zu betrachtenen subjektiven Kriterien fest.
Clarke und Willet messen die Qualität von Suchfunktionen des World Wide Webs anhand des Recalls (\cite{Clarke_Willett_1997}).
Um dies zu ermöglichen wird ein Datensatz generiert, welcher Sucheingaben beinhaltet, sowie alle relevanten Dokumente für eine Sucheingabe.
Bar-Ilan verwendet die gleichen Metriken (\cite{Bar-Ilan_2002}).
Hier wird darüber hinaus auf die Problematik der Messung des Recalls eingegangen.
Es wird erläutert, dass zur Messung des Recalls a-priori bestimmt werden muss, welche Dokumente als relevant für eine gegebene Sucheingabe erachtet werden sollten.
Gordon und Pathak behaupten, dass die Bestimmung der Relevanz lediglich dem Nutzer mit dem Bedürfnis nach der Information zu überlassen ist (\cite{Gordon_Pathak_1999}).
Voorhees und Harman behaupten, dass die Bestimmung der Relevanz durch ein Experten-Panel durchgeführt werden sollte (\cite{Voorhees_Harman_2001}).
Sirotkin betrachtet verschiedene Ansätze zur Messung der Performance von Suchfunktionen (\cite{Sirotkin_2012}).
Neben den bereits genannten Metriken von Precision und Recall werden andere Metriken zur Messung der Performance einer Suchfunktion genannt, wie Mean Reciprocal Rank und Maximal Marginal Relevance.\\

\chapter{Definition von Anwendungsfällen}
\label{chap:anwendungsfaelle}
In diesem Kapitel werden Anwendungsfälle ausgewählt, für welche später Lösungsansätze entwickelt werden.
Die Anwendungsfälle beschreiben die Situationen, in denen Softwareentwickler die Suchfunktionen von Wissensdatenbanken verwenden könnten.
Zur Identifikation von Anwendungsfällen wurde zunächst Literatur herangezogen.
Die Literatur ist bereits unter den verwandten Arbeiten aufgeführt.
In den verwandten Arbeiten wurden Arbeiten genannt, welche ähnliche Probleme lösen sollen.
Diese fokussieren sich vorallem auf die \textit{Feature Location}, \textit{Bug Localization} und die Traceability zwischen Code und anderen Artefakten.
Mit anderen Worten: Die Suchfunktion soll eine Brücke zwischen Code und Dokumentation herstellen.
Neben den beiden Anwendungsfällen Feature Location und Bug Localization aus der Literatur, konnten noch weitere Anwendungsfälle ermittelt werden.
So ist das Onboarding im Projekt ein Anwendungsfäll für die Verwendung der Suchfunktion einer Wissensdatenbank.
Ein weiterer Anwendungsfall ist das Auffinden von Informationen über das Projektmanagement.\\

Die folgenden Kapitel beschreiben die genannten Anwendungsfälle und erläutern, warum sie als Anwendungsfälle ausgesucht wurden.
Bei dem Vergleich zwischen Suchfunktionen werden die Anwendungsfälle herangezogen, um Sucheingaben und erwarteten Dokument zu generieren.
Anhand dieser wird die Studie die Precision der Suchfunktionen ermitteln.

\section{Onboarding im Projekt}
Wenn ein neuer Softwareentwickler in einem Softwareprojekt startet, dann muss er sich zunächst einmal mit dem Projekt vertraut machen.
Das bedeutet, dass er verstehen muss, was das Projekt eigentlich ist.
Er muss verstehen, was das eigentliche Problem des Kunden ist.
Außerdem muss er verstehen wie die Software dieses Problem löst.\\

Für das Onboarding im Projekt muss der Softwareentwickler sehr allgemeine Informationen über das Projekt finden können.
Er könnte Dinge suchen, wie einen Projektüberblick oder ein Glossar.
Neben diesen allgemeinen Informationen muss sich der neue Softwareentwickler mit dem Code vertraut machen.
Er muss verstehen, welche Technologien verwendet werden, welche Best-Practices, Code-Styles, Guidelines, Prozesse und Quality Gates eingehalten werden müssen.
Und er muss verstehen, wie die Software lokal oder in einer Testumgebung ausgeführt werden kann.\\

Die Leitung eines Projektes wünscht sich eine möglichst schnelle Einarbeitung von neuen Mitarbeitern.
Dazu können bereits etablierte Mitarbeiter herangezogen werden, welche den neuen Mitarbeiter bei der Einarbeitung unterstützen.
Der Nachteil besteht darin, dass hierdurch Kapazitäten gebunden werden, welche für die aktive Entwicklung der Software benötigt werden.
Daher kann das Zurückgreifen auf eine Wissensdatenbank durch den neuen Mitarbeiter sinnvoll sein.
Dazu kann der neue Mitarbeiter das Inhaltsverzeichnis der Wissensdatenbank verwenden, wenn es vorhanden ist.
Dieses hilft dem Mitarbeiter dabei Informationen zu bestimmten Themenbereichen zu finden.
Es grenzt die Antwort auf eine Frage auf einen bestimmten Bereich ein, so wie es auch eine Suchfunktion tut.
Nichtsdestotrotz bietet eine Suchfunktion die Möglichkeit dieses Inhaltsverzeichnis automatisch auf Basis einer Sucheingabe zu durchlaufen.
Damit findet der neue Mitarbeiter schneller die Informationen, die er sucht.
Darüber hinaus kann die Suchfunktion spezifischere Ergebnisse liefern.
So kann die Suchfunktion dem Nutzer bereits die relevantesten Stellen in den relevantesten Dokumenten liefern.
Und die Suchfunktion kann dem Nutzer die Informationen aus diesen relevantesten Stellen zusammenfassen.
Voraussetzung hierbei ist eine Suchfunktion, welche gut darin ist die relevantesten Dokumente und die relevantesten Stellen aus diesen Dokumenten zu extrahieren.
Im weiteren Verlauf wird erläutert, welche Ansätze es gibt, um diese Voraussetzung zu erfüllen, und wie aus den Dokumenten die relevantesten Stellen mithilfe von Passage Retrieval ermittelt werden können.
Außerdem wird erläutert, wie Retrieval Augmented Generation die gefundenen Informationen aufbereitet.

\section{Implementierung nach Spezifikation}
Bei der Implementierung von neuen Anforderungen ist es wichtig, dass sich der Softwareentwickler an die Spezifikation hält.
Nur so bekommt der Kunde die Software, die er sich gewünscht hat.
Dazu sollte der Softwareentwickler alle relevanten Dokumente finden, die zu der Spezifikation dazugehören.
Zuerst sollte er die Feature-Spezifikation selbst finden können.
Er sollte die Dokumentation der damit einhergehendenden Prozesse finden, und auch die Domänenobjekte, welche für die Implementierung relevant sind.
Er sollte Diagramme finden können, welche zu dem Anwendungsfall gehören, und auch die weiteren Dokumente, welche den Kontext der Anforderung erläutern.\\

Auch nachdem ein Feature nach Spezifikation umgesetzt wurde, ist es weiterhin wichtig die Spezifikation einfach leicht zu können.
Der Abgleich mit einer Spezifikation ist notwendig, um Testfälle zu schreiben, und zu prüfen, ob das Verhalten der Anwendung korrekt ist.
Es gibt Fehler, welche erkennbar sind, ohne dafür die Spezifikation heranzuziehen.
Eine Anwendung, welche einfriert oder abstürzt ist ein Beispiel für einen solchen Fehler.
Dennoch können Fehler auch domänen- und kontextspezifisch sein, wodurch die Spezifikation als Referenz notwendig ist.
Die Spezifikation zum Abgleich mit dem Verhalten der Software heranzuziehen ist besonders in Randfällen hilfreich.
Angenommen ein Online-Shop bietet einen kostenlosen Versand ab einem Mindestbestellwert an.
Sobald der Preis den Wert von 25\$ übersteigt ist der Versand gratis.
Der Versand kostet im Normalfall 5\$.
Nun muss definiert werden, ob der Versand in dem Mindestbestelltwert miteinbezogen wird oder nicht.
Wird er miteinbezogen, dann sorgt das dafür, dass ein Produkt, welches 20\$ kostet einen kostenlosen Versand hat, weil der Mindestbestellwert zuzüglich dem Versandpreis 25\$ beträgt.
Ob dieses Verhalten gewünscht ist oder nicht, muss in der Spezifikation festgehalten werden.
Wenn nun ein Bugticket dieser Art bei einem Softwareentwickler landet, dann muss er, wie auch bei der Implementierung, alle relevanten Dokumente für die weitere Entwicklung heranziehen.//

Neben dem Auffinden der relevanten Dokumente aus der Wissensdatenbank kann die Suche einen weiteren Mehrwert bieten, wenn sie auch gleich die relevante Stelle im Code findet, welche angepasst werden muss, um das Feature zu implementieren oder zu erweitern, oder um einen Bug zu beheben.
\textit{Feature Location} ist das Vorgehen, den Einstiegspunkt im Code zu finden, welcher für ein Feature verantwortlich ist.
Wenn eine Software das Feature besitzt, Mails zu verschicken, dann könnte der Einstiegspunkt für dieses Feature eine Schnittstelle sein, welche den Inhalt und Empfänger der Email entgegennimmt.
Der Inhalt wird anschließend in ein Mail-Template eingefügt, und die Mail wird verschickt.
Wenn die Suchfunktion nun Feature Location ermöglicht, dann wird damit die Erweiterung des Features vereinfacht, weil der Software-Entwickler die anzupassende Stelle nicht manuell suchen muss.\\

Analog hierzu ist \textit{Bug Localization} ein Verfahren, um die Quelle von Fehlern im Code einzugrenzen.

\section{Bug Localization}
Neben der Implementierung neuer Features ist es Aufgabe eines Softwareentwicklers bestehende Fehler in der Software zu korrigieren.
Um einen Fehler zu korrigieren, ist es sinnvoll die Spezifikation heranzuziehen, um herauszufinden, wie sich die Software ordnungsgemäß verhalten sollte.
% TODO: Neu schreiben

Wenn der Softwareentwickler nun ein Fehlerticket erhält, dann muss er die entsprechende Spezifikation zu diesem Fehlerticket finden können.
Und idealerweise wird ihm durch die Suche sogar gleich die betroffene Stelle im Code angezeigt.
Das ist Bug Localization.

\section{Informationen über das Projektmanagement finden}
Softwareentwickler sollten darüber Bescheid wissen, was in dem Projekt, in dem sie arbeiten, vor sich geht.
Sie sollten über organisatorische Dinge Bescheid wissen, wie die Teamaufteilung und die Zuständigkeiten in den einzelnen Teams und in dem gesamten Projekt.
Das ist wichtig für eine gute Kommunikation und entsprechenden Wissensaustausch zwischen Kollegen.
Außerdem sollten Softwareentwickler darüber Bescheid wissen, wann das nächste Release der Software ansteht.
Diesen Termin müssen sie bei der Planung ihrer Arbeit berücksichtigen, damit sie auch rechtzeitig alle relevanten Features implementiert haben und alle kritischen Fehler behoben haben.
Neben der Einhaltung von Terminen ist die Planung wichtig für die Motivation der Softwareentwickler.
Denn Motivation entsteht durch die Wahrnehmung, sich einem Ziel zu nähern (Quelle: Flow - mihaly csikszentmihaly).
Aus diesem Grund müssen die Softwareentwickler auch die Vision der Software verstehen und auch die übergreifende Vision des Unternehmens.
Sie müssen verstehen, warum die Arbeit, welche sie erledigen so wichtig ist.
Und sie müssen verstehen, für wen sie diese Arbeit tun.
Auch diese Faktoren sind wichtig für eine hohe Motivation bei der Arbeit.
Daher ist es so wichtig diese Informationen einfach zugänglich zu machen, also einfach auffindbar zu machen.

In \myRef{Kapitel}{chap:evaluationsmethoden} wird aufgezeigt, wie die Anforderungen der Anwendungsfälle quantifizierbar gemacht werden können.
Das wird dabei helfen nachzuvollziehen, inwieweit die Anforderungen der Anwendungsfälle erreicht wurden, welche genannt wurden.

\chapter{Evaluationsmethoden und -Kriterien}

\section{Features}

\subsection*{Relevanz}
Eine Suchfunktion muss nicht nur Ergebnisse für eine Such-Query finden, sondern diese auch darstellen. Bei der Darstellung der Ergebnisse ist die Reihenfolge besonders wichtig. Manchmal hat der Nutzer keine genaue Vorstellung davon, was genau er sucht. Er möchte Informationen zu einem Thema finden, aber ohne dabei bestimmte Seiten oder Inhalte im Hinterkopf zu haben. Nichtsdestotrotz möchte der Nutzer mit minimalem Aufwand so viele wichtige Informationen, wie möglich erhalten. Die Sortierung von Ergebnissen der Suche soll genau das ermöglichen.
Nun ist zu klären, wie die Bestimmung der Relevanz von Suchergebnissen überhaupt funktioniert. Es muss bestimmt werden, anhand welcher Kriterien die Suchfunktion die Relevanz bestimmt. Für diese Kriterien gibt es wiederum unterschiedliche Quellen. Eine Quelle für Kriterien der Relevanz ist das Nutzerverhalten. Seiten, die besonders häufig besucht werden und auf denen lange verweilt wird, oder welche sehr häufig bearbeitet werden, sind gute Kandidaten für eine hohe Relevanz. Seiten, welche bei Nutzern praktisch in Vergessenheit geraten sind, scheinen nicht so wichtig zu sein, sodass die Relevanz für die Suche niedrig ist.   

\subsection*{Clustering}
Thematisch ähnliche Dokumente tendieren dazu relevant für die gleiche Suche zu sein. Aus diesem Grund ist es oftmals sinnvoll ähnliche Dokumente zu clustern, also zu gruppieren, und sie dem Nutzer als ein solches Cluster darzustellen. Alternativ kann auch ein Suchergebnis dargestellt werden und ähnliche Suchergebnisse werden angezeigt als „Ähnliche Suchergebnisse“.

\subsection*{Autovervollständigung}
Autovervollständigung bedeutet, dass während der Eingabe von Suchbegriffen bereits Ergänzungen für die gesamte Eingabe vorgeschlagen werden. Das ist vor allem ein hilfreiches Feature in Fällen, wo der Nutzer nur eine grobe Ahnung davon hat, was er eigentlich sucht. Mithilfe von Vorschlägen durch die Suchmaschine selbst kann der Nutzer seine Eingabe spezifischer fassen als er es ohne Vorschläge schaffen würde, sodass die Suchergebnisse besser zu dem passen, was er eigentlich sucht.

\subsection*{Best-Prefix-Highlighting}
Das ähnlichste Ergebnis zur Eingabe des Nutzers soll gehighlightet werden.

\subsection*{Synonyme}
Synonyme bei Suche mitbeachten.

\section{Qualitätskriterien}

\subsection*{Interaktivität}
Klassische Suchfunktionen sind so konzipiert, dass der Nutzer einen Suchbegriff eingibt, anschließend auf einen Submit-Button klickt und erst dann Suchergebnisse angezeigt bekommt. Der Nachteil dieser Vorgehensweise ist, dass der Nutzer erst spät Feedback darüber bekommt, ob er mit den eingegebenen Suchbegriffen auch Ergebnisse bekommt, die für ihn gerade sinnvoll sind. Das Qualitätskriterium Interaktivität gibt an, wie kurz die Zeit zwischen dem Eingeben von Suchbegriffen und der Anzeige von Ergebnissen ist. Eine Suchfunktion muss nicht unbedingt einen Submit-Button verwenden. Die Suche kann auch mit jedem ausgeschriebenen Keyword die Suche starten und Ergebnisse dafür anzeigen, oder auch nach jedem eingegebenen Buchstaben. Mit einer kürzeren Zeit zwischen Eingabe und Ergebnis, also einer höheren Interaktivität, weiß der Nutzer schneller, ob seine Eingaben sinnvoll sind, und bekommt damit schneller gute Suchergebnisse.

\subsection*{Geschwindigkeit}
Die Geschwindigkeit einer Eingabe ist die Zeit zwischen dem Beginn der Suche und der Anzeige der ersten Antwort. Die Geschwindigkeit hat nichts mit der Interaktivität zu tun. 

\subsection*{Darstellung}
Wie gut kann man anhand der Zusammenfassung der Ergebnisse den Inhalt bestimmen.

\subsection*{Benutzerfreundlichkeit}
Das System findet auch Ergebnisse für ähnliche Wörter. So sind bspw. Tippfehler kein Problem. Fuzzy-Search.

\section{Precision, Recall und F-Maß}
Zur Evaluation werden die Werte Precision und Recall gemessen. Der Precision-Wert ist das Verhältnis zwischen allen gefundenen Dokumenten und den gefundenen Dokumenten, die tatsächlich relevant sind. Der Recall-Wert gibt an, wie viele von den tatsächlich relevanten Dokumenten auch gefunden wurden. Es ist schwierig beide Werte zu optimieren, da der Precision-Wert versucht die Anzahl der gefundenen Dokumente einzugrenzen und der Recall-Wert versucht die Anzahl der gefundenen Dokumente zu erweitern. Das F-Maß fasst beide Werte zu einem neuen Wert zusammen.

\section{Befragungen}
TODO
\chapter{Theoretischer Hintergrund}
Die Architektur von Suchfunktionen besteht aus drei Komponenten.
Aus dem Crawling, der Indexierung und der Suchfunktion selbst.
Das Crawling ist zuständig für das Finden von Websites.
Die Indizierung ist zuständig für das optimale Speichern der Informationen der Websites, und die Suche ist zuständig für das Verstehen der Nutzeranfrage und die Abfrage der relevantesten Informationen aus dem Index, sowie dessen Verarbeitung und Darstellung.
Im Folgenden wird der Begriff Dokument verwendet, um die Dateien zu beschreiben, welche durch einen Crawler gesucht und durch den Index verarbeitet werden.
Unter Dokumenten können hierbei auch eine Website verstanden, welche durch einen Webcrawler durchsucht werden.
Es wird zunächst die Funktionsweise des Crawlings erläutert.
Danach werden zwei Möglichkeiten der Indizierung beschrieben.
Zuerst die Volltext-Indizierung, dann die Vektor-Indizierung.
Zuletzt werden die verschiedenen Arten von Suchalgorithmen erläutert.

\section{Crawling}
Bei der Implementierung eines Systems für eine Suchfunktion benötigt das System zunächst einen Datensatz von Dokumenten, welche überhaupt grundsätzlich über die Suchfunktion gefunden werden können.
Dieser wird mithilfe eines Crawlers aufgebaut.
Ein Crawler ist ein Algorithmus, welcher ein Dokument als Startpunkt bekommt, und anhand dessen neue Dokumente findet.
Der Algorithmus analysiert dazu das Dokument auf Links, welchen der Algorithmus anschließend folgt.
Die neuen Dokumente werden durch den Index verarbeitet und wiederum auf neue Links analysiert.
Dieses Verfahren kann beliebig lange und beliebig rekursiv durchlaufen werden, um den Index zu erweitern.
Neben dem Crawling können Indizes befüllt werden, indem eine Liste von Dokumenten übergeben werden, welche dem Index hinzugefügt werden sollen.

\section{Indizierung}
Bei der Indexierung werden Wörter und Tokens mit Dokumenten assoziiert.
Dazu werden Wörter aus gecrawlten Dokumenten extrahiert und in die Datenbank geschrieben.
Den Dokumenten werden IDs zugeordnet und diese IDs werden den Wörtern zugeordnet.
Dem Index werden nun weitere Informationen hinzugefügt, wie die Wortfrequenz.
Die Wortfrequenz gibt an, wie oft ein Wort in einem Dokument vorkommt.
Es wird auch gespeichert, an welchen Stelle des Dokuments das Wort vorkommt, und auch in wie vielen Dokumenten ein Wort vorkommt.
Diese Art der Indizierung wird auch als invertierter Index bezeichnet, weil den Wörtern bzw.
Tokens die Dokumente zugeordnet werden, und nicht umgekehrt.
Wenn der Nutzer nun ein Keyword in die Suche eingibt, dann sind diese Keywords oft bereits im Index vorhanden, sodass die Dokumente, welche diese Keywords beinhalten einfach dem Index entnommen werden können.
Aufgabe der Suche ist es anschließend die Ergebnisse aufzubereiten.

\subsection{Volltext-Indizierung}
Bei der Indizierung der Wörter besteht die Problematik, dass gleiche Wörter in unterschiedlichen Formen existieren können.
So stammen \textit{Heizung} und \textit{heizen} beide von dem gleichen Wortstamm \textit{heiz} ab.
Um bei der Indizierung Speicherplatz zu sparen, können Wörter auf diesen Wortstamm reduziert werden, damit sie als ein einziges Wort betrachtet werden können.
Die Bildung des Wortstamms wird auch als Stemming bezeichnet.
Beim Stemming kann es jedoch zu Overstemming und Understemming kommen.
Overstemming bedeutet, dass zwei Wörter, die eigentlich nichts miteinander zu tun haben, also nicht semantisch gleich sind, den gleichen Wortstamm besitzen und als ein Wort betrachtet werden.
Ein Beispiel hierfür sind die Wörter \textit{Wand} und \textit{wandere}, wie in \textit{ich wandere}.
Beide besitzen den Wortstamm \textit{wand} und werden entsprechend als ein Wort betrachtet.
Understemming bedeutet, dass zwei Wörter, die eigentlich etwas miteinander zu tun haben, also semantisch gleich sind, nicht den gleichen Wortstamm besitzen und dadurch als zwei verschiedene Wörter betrachtet werden.
Ein Beispiel hierfür sind die Wörter \textit{absorbieren} und \textit{Absorption}, welche die Wortstämme \textit{absorb} und \textit{absorp} besitzen.
Es gibt Techniken zur Vermeidung solcher Probleme, wie der Einsatz vollständiger morphologischer Analysekomponenten.
Hierauf soll aber nicht weiter eingegangen werden.\\

Zur Implementierung eines Volltext-Index werden Dokumente und Wörter in einer m x n Matrix angeordnet, wobei m die Anzahl der Dokumente ist und n die Anzahl der Wörter.
Die Werte in dieser Matrix werden anhand einer ausgewählten Metrik bestimmt.
Eine oft verwendete Metrik ist das Verhältnis der Wortfrequenz mit der invertierten Dokumentfrequenz.
Die Wortfrequenz gibt dabei an, wie häufig das in dem Dokument vorkommt.
Die Dokumentfrequenz gibt an, in wie vielen Dokumente ein Wort vorkommt.
Das Verhältnis gibt damit an, wie charakteristisch das Wort für ein bestimmtes Dokument ist.
Wenn das Wort charakteristisch für ein Dokument ist, dann kommt es in diesem Dokument häufig vor, aber in anderen Dokumenten nur selten.
Wenn das Wort nicht charakteristisch für das Dokument ist, dann kommt es in anderen Dokumenten genauso häufig oder häufiger vor als in diesem Dokument.
Die Metrik wird auch als tf-idf-Wert bezeichnet.
TODO: Rechnung besser verstehen\\

\subsection{Vektor-Indizierung}
Neben einer Volltext-Indizierung können Dokumente in Form von Vektoren indiziert werden.
Dazu werden Dokumente zunächst, wie auch bei der Volltext-Indizierung gecrawlt.
Anschließend durchlaufen die Inhalte der Dokumente ein Preprocessing.
Dieses kann je nach Implementierung variieren.
Das Kapitel \textit{Einspielen der Daten in Weaviate} beschreibt, wie in der hier aufgeführten Implementierung das Preprocessing durchgeführt wird.
Das Preprocessing hat den Zweck die Daten an das Schema der Datenbank anzupassen und die Qualität der Daten zu erhöhen.
Außerdem sorgt es für eine kürzere Indizierungszeit.\\

Nach dem Preprocessing werden durch einen Transformer für die Inhalte der Dokumente Vektoren berechnet.
Ein Transformer wird mithilfe von Trainingsdaten darauf trainiert, Vektoren für Wörter zu generieren.
Das trainierte Modell wird nach dem Preprocessing durchlaufen.
Zuletzt werden die Daten in einer Vektordatenbank gespeichert.
Eine Vektordatenbank speichert Daten, wie eine dokumentenbasierte Datenbank.
Dort werden nun sowohl die rohen Daten als auch die Vektoren gespeichert, welche von dem Transformer berechnet wurden.
Die Vektoren haben den Vorteil, dass die Daten in der Datenbank nicht linear gespeichert sind.
Sie sind in einem n-dimensionalen Raum gespeichert, mit dessen Hilfe die semantische Nähe zwischen Dokumenten ausgedrückt werden kann.
Das funktioniert auf die gleiche Weise, wie bereits in dem Kapitel \textit{Semantic Search} beschrieben.

\section{Suchalgorithmen}
Im Folgenden werden zuerst gängige Suchalgorithmen erklärt, welche in nahezu jeder Suchfunktion zu finden sind.
Danach wird im speziellen auf die strukturierte Suche und die semantische Suche eingegangen.
Die semantische Suche wird später bei der Implementierung einer Suchfunktion verwendet werden.

\subsection{Die gängigen Suchalgorithmen}
In dem Kapitel \textit{Evaluationsmethoden und -Kriterien} wurde bereits beschrieben, dass eine Suche das Qualitätskriterium der Konformität erfüllen sollte.
Das bedeutet, dass eine Suchfunktion die gängigen Arten von Suchanfragen unterstützen muss.
Dieses Kapitel soll die Arten von Suchanfragen vorstellen.
Es ist wichtig diese zu kennen, um das Qualitätskriterium später bei der Implementierung erfüllen zu können.
Es gibt die gängigen Suchalgorithmen Keyword Search, Phrase Search, Boolean Search, Field Search.\\

Eine Keyword Search durchsucht Dokumente nach der Sucheingabe des Nutzers.
Die Eingabe wird dabei nicht als ganzes betrachtet, sondern jedes Wort einzeln.
Für jedes Keyword werden Dokumente als Ergebnis angezeigt, wenn dieses in dem Dokument vorhanden ist.
Wenn ein Dokument mehrere der Keywords beinhaltet, wird dessen Relevanz höher eingeschätzt als für Dokumente, welche weniger Keywords enthalten.
Dokumente mit höherer Relevanz werden weiter oben in der Ergebnisliste angezeigt.\\

Eine Phrase Search ist die Suche nach Textausschnitten in Dokumenten.
Hier werden nicht mehrere Keywords einzeln betrachtet, sondern die gesamte Eingabe in das Suchfeld als eine Einheit.
Es reicht also nicht mehr aus, dass ein Dokument eines der Wörter enthält.
Es muss die gesamte Sucheingabe als ein String enthalten sein.\\

Die Boolean Search bietet die Möglichkeit einen boolschen Ausdruck als Sucheingabe zu machen.
Ein Beispiel dafür ist die Sucheingabe \textit{Dokumentation AND Angular}.
Die Sucheingabe bedeutet, dass die Suchfunktion nur Dokumente als Ergebnis darstellen soll, welche beide Keywords Dokumentation und Angular enthalten.
Die Boolean Search kann auch eine Phrase, wie bei der Phrase Search, beinhalten: \textit{"Dokumentation von Software" AND Angular}.
In diesem Beispiel werden nur Dokumente als Ergebnis nur angezeigt, wenn sie den gesamten String \textit{Dokumentation von Software} enthalten, sowie das Keyword \textit{Angular}.
Bei einer Boolean Search können die boolschen Operatoren \textit{AND}, \textit{OR}, \textit{NOT} beliebig kombiniert werden.\\

Eine Field Search sucht Dokumente anhand von Attributen.
Der Nutzer kann diese Attribute auswählen.
Wenn der Nutzer beispielsweise ein Dokument sucht, welches am 01.01.2005 erstellt wurde, dann kann die Eingabe der Suche so aussehen: \textit{erstelldatum: 01.01.2005}.
Es können beliebig viele Attribute verwendet werden, um die Suche einzugrenzen.
Neben der Verwendung der Attribute für die Suche selbst, können die Attribute komplementär zu einer anderen Art von Suche verwendet werden.
So kann eine Suchfunktion Buttons bereitstellen, über welche Filter festgelegt werden.
Nun kann eine Keyword Search durchgeführt werden, aber die gefundenen Dokumente werden mithilfe der Filter weiter eingeschränkt.
Neben diesen gängigen Suchalgorithmen gibt es weitere Suchalgorithmen, wie die strukturierte Suche und die semantische Suche.

\subsection{Die strukturierte Suche}
Eine klassische Suchfunktion verwendet Keywords, um relevante Dokumente für eine Sucheingabe zu ermitteln.
Der Vorteil dieser Art von Suche ist, dass die technische Struktur, in der die Daten vorliegen und gespeichert sind, nicht bekannt sein müssen.
Der Nachteil ist auf der anderen Seite, dass die Suchergebnisse unpräzise sein können.
Wenn beispielsweise der Suchbegriff Kamera eingegeben wird, dann werden Kameras als Ergebnisse zurückgegeben.
Soll die Kamera nun bestimmte Eigenschaften besitzen, dann müssen diese Eigenschaften ebenfalls als Keywords angegeben werden.
Nun wird aber nicht die Suche auf Ergebnisse eingegrenzt, bei denen eine Kamera diese bestimmten Eigenschaften besitzt.
Stattdessen werden Suchergebnisse angezeigt, bei denen einige dieser Keywords vorkommen.
Demgegenüber stehen Datenbankabfragen, beispielsweise mithilfe von SQL.
Bei einer Datenbankabfrage können Objekte abgefragt werden, dessen Eigenschaften ganz bestimmte Werte haben.
Die Ergebnisse, die eine solche Abfrage zurückgibt, sind dabei vollkommen genau.
Es werden keine Objekte zurückgegeben, welche diese Kriterien nicht erfüllen.
Voraussetzung für eine solche Suchabfrage ist allerdings, dass die Struktur der Datenbank a priori bekannt ist.
Dem Nutzer muss der Name der Datenbank, der relevanten Tabellen, sowie der relevanten Properties bekannt sein, damit er das passende SQL für die Datenbankabfrage schreiben kann.
Strukturierte Suchen sollen die Vorteile beider Vorgehensweisen kombinieren.
Die Struktur der Daten soll a priori nicht bekannt sein müssen, aber trotzdem sollen die Suchergebnisse vollkommen präzise sein.
TODO: Wie wird das ermöglicht?

\subsection{Die Semantische Suche}
Unter einer semantischen Suche wird im Allgemeinen verstanden, dass die Suche nicht nur eine syntaktische Suche einer Zeichenkette ist, sondern auch Techniken, wie technische Analysen verwendet, um die Bedeutung der Sucheingabe nachzuvollziehen.\cite{Dengel_2012}
Eine semantische Suche hat den Zweck, die Ähnlichkeit und Beziehungen zwischen Wörtern zu verstehen.
Sie kennt Homonyme, Synoyme und Antonyme von Wörtern.  
So wird durch sie beispielsweise die Ähnlichkeit von den Wörtern \textit{rollout} \textit{deployment} abgebildet, und dass diese Wörter oft im gleichen Kontext verwendet werden.\\

Die technische Umsetzung einer semantischen Suche kann auf verschiedene Arten erfolgen.
Zum einen besteht die Möglichkeit ein explizites semantisches Netz heranzuziehen, und so die Zusammenhänge von Wörtern abzubilden.
Ein semantisches Netz ist eine Menge aus Aussagen in der Form \textit{Subjekt Prädikat Objekt}.
Anhand dieser Aussagen wird ein Graph aus Beziehungen zwischen Wörtern hergestellt.
TODO: Erläuterende Grafik + Beispiel.
Das hat den Vorteil, dass ein solches explizites semantisches Netz von Menschen erstellt wurde, und damit eine hohe Qualität der Daten einhergeht.
Denn an der Erstellung von semantischen Netzen sind mehrere Menschen beteiligt, welche zuerst einen Konsens über die Eigenschaften und Zusammenhänge von Wörtern schaffen müssen.
Der Nachteil dieser Methode ist der große Arbeitsaufwand, welcher mit der Erstellung eines solchen semantischen Netzes einhergeht.
Ein weiterer Nachteil ist die mögliche Unvollständigkeit, welche ein solches semantisches Netz besitzen könnte.
Eine Wissensdatenbank, welche für ein Projekt erstellt wurde, kann Definitionen von Begriffen beinhalten, welche in allgemeinen semantischen Netzen nicht vorhanden sind.
Bei der Umsetzung einer semantischen Suche mithilfe eines semantischen Netzes müsste also zuerst ein allgemeines semantisches Netz herangezogen werden, beispielsweise von DBPedia.
Anschließend müsste dieses semantische Netz um die neuen Definitionen, welche nur im Kontext des Projektes gelten, erweitert werden.\\

Eine weitere Möglichkeit zur Umsetzung einer semantischen Suche ist die Verwendung von Transfomers und Vektordatenbanken.
Um zu verstehen, welche Wörter kontextuell zusammengehören, werden hier die Wörter in einem n-dimensionalen Raum positioniert.
Wörter, die sich sehr ähnlich sind, also im gleichen Kontext verwendet werden, haben in diesem n-dimensionalen Raum eine geringe Distanz zueinander.
Wörter, die sich eher unähnlich sind, wie \textit{"rollout" und "API"}, haben eine größere Distanz.
Der Vorteil einer semantischen Suche ist, dass der genaue Begriff, welcher gesucht wird nicht bekannt ist.
Wenn sich der Nutzer also über ein Thema informieren möchte, mit welchem er nicht gut vertraut ist, dann kann die semantische Suche hilfreich sein.
Denn der Nutzer kann nun einen Begriff eingeben, der zu dem Thema passt, und den er kennt.
Er findet anschließend Dokumente, welche vielleicht nicht genau diesen Begriff beinhalten, aber welche thematisch dennoch ähnlich sind.
Genau dieser Vorteil soll bei der Implementierung später genutzt werden.\\

Um eine semantische Suche zu implementieren, werden die Technologien von Transformern und Vektordatenbanken verwendet.
Ein Transformer erhält als Input eine große Menge an Text und mappt die einzelnen Wörter auf einen Vektor einer beliebigen Länge.
Der Vektor, der am Ende herauskommt, beschreibt die Position des Wortes in dem n-dimensionalen Raum.
Der Vektor beschreibt gewissermaßen, wie stark ein Wort in eine abstrakte Kategorie einzuordnen ist.
Jeder Wert im Vektor entspricht einer Kategorie.
Mithilfe der Vektoren können verschiedene Wörter hinsichtlich ihrer Ähnlichkeit analysiert werden.
Ähnliche Wörter habe eine große räumliche Nähe, während zwei Wörter, die in vollkommen unterschiedlichen Kontexten verwende werden eine sehr große Distanz im Raum besitzen.
Nehmen wir für ein Beispiel einen dreidimensionalen Raum an.
Die X-Achse ist beschriftet mit dem Wort „Tier“, die Y-Achse ist beschriftet mit dem Wort „Computer“ und die Z-Achse ist beschriftet mit dem Wort „Mensch“.
Nun geben wir einem Transformer das Wort „Katze“, und der Transformer berechnet einen dreidimensionalen Vektor, welcher das Wort „Katze“ im Raum positioniert.
Weil eine Katze ein Tier ist, ist der X-Wert des Vektors eins.
Der Wert eins bedeutet, dass das Wort vollständig zu dieser Kategorie gehört.
Da eine Katze überhaupt nichts mit einem Computer zu tun hat, ist der Y-Wert des Vektors 0.\\

Nun ist eine Katze kein Mensch, aber eine Katze ist ein Haustier von Menschen.
Es ist denkbar, dass die Wörter Katze und Mensch oft im gleichen Kontext verwendet werden, sodass der Wert bei 0,3 liegen könnte.
Damit der Transformer einen Vektor berechnen kann, braucht er eine Menge Daten.
Diese Daten erhält er aus vielen Texten.
Werden zwei Wörter oft im gleichen Text genannt oder kommen zwei Wörter in vielen Texten sehr nahe beieinander vor, dann geht der Transformer davon aus, dass die beiden Wörter ähnlich sind, und berechnet ähnliche Vektoren.
Zuvor müssen die Texte allerdings bereitgestellt werden.
Dazu kann beispielsweise das Internet gecrawlt werden.
Die Ergebnisse des Transformers werden in einer Vektordatenbank gespeichert.
Eine Vektordatenbank ist eine Datenbank, welche Vector Embeddings, also ein Objekt als Key und dessen Vektor als Value speichert.
Bei dem Objekt kann es sich um Wörter handeln, dann wird auch von Word Embeddings gesprochen.
Es können aber auch Daten andere Daten, wie Bilder, Videos oder Audio gespeichert werden.
Der Zweck von Vektordatenbanken ist es, Daten nicht einfach linear zu speichern, sondern in einem Raum.
Die Distanz zwischen zwei Einträgen in diesem Raum beschreibt dessen Ähnlichkeit.
Genau diese Informationen nutzen semantische Suchen.
\chapter{Implementierung der neuen Suchfunktion}
\label{chap:implementierung}
In diesem Kapitel wird die Implementierung der neuen Suchfunktion dargestellt.
Die Implementierung verwendet eine Vektordatenbank, welche eine semantische Suche ermöglicht.
Es wird die Vektordatenbank Weaviate verwendet.
% TODO: Weaviate Benchmarks darstellen und damit begründen, dass die semantische Suche besser sein sollte

\section{Aufsetzen der Vektordatenbank Weaviate}
Für das Aufsetzen von Weaviate werden zwei Komponenten benötigt.
Zum einen die Vektordatenbank selbst.
Zum anderen ein Transformer, welcher Text entgegennimmt, und diese in Vektoren umwandelt, sodass diese in der Datenbank gespeichert werden können.
Um die beiden Komponenten aufzusetzen wird Docker verwendet.
Das entsprechende docker-compose.yml File ist im Anhang zu finden.
Hier werden die beiden Komponenten definiert.
Unter \textit{t2v-transformers} wird der Transformer konfiguriert.
Es wird das Image eines bereits vortrainierten Transformers verwendet.
Unter \textit{weaviate} wird die Datenbank konfiguriert.
Hier wird mithilfe von den Environment-Variablen \textit{$DEFAULT\_VECTORIZER\_MODULE$}, \textit{$ENABLE\_MODULES$} und \textit{$TRANSFORMERS\_INFERENCE\_API$} konfiguriert, welcher Transformer verwendet werden soll.
So wird konfiguriert, dass der Transformer der anderen Komponente verwendet werden soll.\\

Neben den beiden Docker-Containern für die semantische Suche wird eine Spring Boot Anwendung in einem Docker Container mitdeployt.
Die Spring Boot Anwendung verwendet Kotlin und Gradle.
Mithilfe der Dependency \textit{io.weaviate:client:4.0.1} wird die Client API von Weaviate verwendet.
So wird nun eine Verbindung zu der Vektordatenbank aufgebaut.
Diese beinhaltet zu diesem Zeitpunkt noch keine Daten.
Bevor die Daten in die Datenbank eingespielt werden, muss das Schema der Daten angegeben werden.
Der entsprechende Code ist ebenfalls im Anhang zu finden.
Es wird eine Klasse \textit{Document} definiert.
Diese Klasse beinhaltet die Properties \textit{documentUrl}, \textit{h1}, \textit{h2} und \textit{p}.
Es werden in dieser Klasse also die URL des Dokuments, sowie die Inhalte aller h1-, h2 und p-Tags gespeichert.
Außerdem wird als Vectorizer \textit{text2vec-transformers} konfiguriert.\\

\section{Einspielen der Daten in Weaviate}
\label{chap:einspielen-der-daten-in-weaviate}
Um nun die Daten aus Confluence in Weaviate einzuspielen, werden zuerst die Daten aus Confluence exportiert.
Beim Export von Confluence-Seiten werden HTML-Dateien generiert.
Es werden keine CSS-, JavaScript- oder Bilddateien generiert.
Die HTML-Dateien werden preprocessed, um dem zuvor definierten Schema zu entsprechen.
Es werden zuerst mithilfe eines Abstract Syntax Trees alle Inhalte von h1-, h2- und p-Tags herausgefiltert.
Anschließend werden Punctuations und Stopwords entfernt und die Inhalte druchlaufen einen Tokenizer und einen Stemmer.
Punctuations sind Zeichen, wie die folgenden: \textit{.,;:}.
Stopwords sind Wörter, welche für einen Leser notwendig sind, aber für die Verarbeitung durch einen Algorithmus als unwichtig erachtet werden (\cite{Sarica_Luo_2021}).
Beispiele für Stopwords sind \textit{aber, denn, der}.
Ein Tokenizer trennt einen Text in einzelne Wörter auf.
Aus einem Text, wie \textit{das deployment erfolgt durch ein bash-skript} wird also ein Array, welches folgendermaßen aussieht:\\

\(["das", "deployment", "erfolgt", "durch", "ein", "bash", "skript"]\).\\

Ein Stemmer bestimmt den Wortstamm für die einzelnen Wörter auf Basis von Grammatikregeln.
Die Software verwendet den Porter-Stemmer-Algorithmus.
Aus dem Wort \textit{deployment} wird dadurch beispielsweise \textit{deploy}.
Duplikate von Wörtern werden anschließend verworfen.
Alle Inhalte, welche in h1-Tags stehen, werden konkateniert und einem Datenfeld der Vektordatenbank gespeichert.
Das gleiche gilt für die übrigen Tags.
Nachdem die Inhalte das Preprocessing durchlaufen haben, werden sie in die Datenbank eingespielt.\\

In \myRef{Kapitel}{chap:vektorindizes} wurde bereits von Performancegründen gesprochen, aus denen das Preprocessing durchgeführt wird.
Da nun die einzelnen Schritte des Preprocessings erklärt wurden, kann auch erklärt werden, warum das Preprocessing die Performance beim indizieren erhöht.
Zuallererst werden viele Wörter gänzlich verworfen, weil sie Stopwords sind.
Durch den Stemmer werden anschließend ähnliche Wörter zusammengruppiert.
Die Wörter \textit{Heizung} und \textit{heizen} stammen beide vom gleichen Wortstamm \textit{heiz}.
Das bedeutet, dass aus zwei verschiedenen Wörtern, welche beide indiziert werden müssen, ein einziges Wort gemacht wird.
Denn am Ende werden Duplikate, wie bereits erwähnt, verworfen.
Die Anzahl der Wörter, welche indiziert werden müssen, wird dadurch reduziert, und damit auch die Last auf dem Transformer, welche die Vektoren für die Wörter berechnen muss.
\chapter{Vergleich der Suchfunktionen}
\label{chap:vergleich-der-suchfunktionen}

Zum Vergleich der beiden Suchfunktionen wird im Folgenden ein fiktionales Szenario dargestellt.
Das Szenario beschreibt den Ablauf des Onboardings eines neuen Mitarbeiters.
Dabei werden die expliziten Sucheingaben beschrieben, die der neue Mitarbeiter in der Wissensdatenbank macht, um sich einzuarbeiten.
Die Ergebnisse dieser Sucheingaben werden im nächstes Schritt dargestellt und ausgewertet.
Die zugrundeliegenden Daten sind im Anhang zu finden.
Zuletzt wird auf die Validität der Daten eingegangen.
Es wird der Versuchsaufbau diskutiert.

\section{Fiktionales Szenario: Onboarding eines Mitarbeiters}

Für den Vergleich von Suchfunktionen müssen auf Grundlage der Anwendungsfälle realistische Suchanfragen entwickelt werden.
Zu diesem Zweck sei angenommen, dass ein neuer Softwareentwickler in einem bestehenden Softwareprojekt eingearbeitet wird.
Es wird also der Anwendungsfall des Onboardings betrachtet.\\

Zuerst wird dem Softwareentwickler aufgetragen, sich das \textit{Getting Started} im Confluence durchzulesen.
Er macht also die Sucheingabe \textbf{Getting Started}, und erwartet ein Dokument mit ebendieser Überschrift.

TODO: Sucheingabe weiter beschreiben

Um sich mit der Software vertraut zu machen, wird dem neuen Softwareentwickler aufgetragen, die Anwendung bei sich lokal zu starten.
Nachdem er sie gestartet hat, soll er sich mit der Funktionalität der Software vertraut machen.
Um die Anwendung lokal zu starten, sucht der Softwareentwickler nach einer \textbf{Installationsanleitung}.
Er gibt also genau dies als Sucheingabe ein und erwartet als Ergebnis ein Dokument, welches beschreibt, wie die Software installiert wird.
TODO: Sucheingabe weiter beschreiben\\

Nachdem die Software installiert ist, startet er die Anwendung.
Dazu muss er sich anmelden.
Er sucht also nach \textbf{Testdaten}, welche die Anmeldedaten enthalten.
Nachdem er sich angemeldet hat, möchte er sich eine Übersicht über die Funktionen der Software verschaffen.
Dazu sucht er nach den \textbf{Use-Cases} der Anwendung.
Nachdem er die Use-Cases gefunden hat, probiert er mehrere davon aus.\\

Einer der Use-Cases ist die Versendung von Korrespondenzen am Ende eines Workflows.
Um diesen Use-Case genauer nachvollziehen zu können gibt er den Suchbegriff \textbf{Korrespondenz} ein.

Ein weiterer Use-Case ist die Erstellung eines Auftrags.
Die genauen die Details der Software dürfen nicht geschildert werden.
Daher kann nicht weiter darauf eingegangen werden, was ein Auftrag genau ist.
Er macht die Sucheingabe \textbf{Auftrag erstellen}.\\

Nachdem ein Auftrag erstellt wurde, lassen sich die Daten eines Auftrages anpassen.
Unter anderem lässt sich die Email-Adresse ändern, an welche die Korrespondenz am Ende des oben genannten Workflows geschickt wird.

Der Softwareentwickler möchte herausfinden, wie er diese ändern kann.
Dazu macht er die Sucheingabe \textbf{Empfänger von Korrespondenz ändern}.\\

Der Softwareentwickler möchte sich weiterhin noch technischer mit der Software auseinandersetzen, und möchte sich daher das genaue Datenmodell eines Auftrags ansehen.
Er macht also die Sucheingabe \textbf{Datenmodell Auftrag}.\\

Nachdem sich der Softwareentwickler eine Weile mit der Software auseinandergesetzt hat, wird ihm mitgeteilt, dass Mitarbeiter des Projektes in Teams eingeteilt sind.
Dabei gibt es drei Entwicklungsteams: Team Blau, Team Rot, Team Gold.
Ihm wird erklärt, dass er von nun an Teil von Team Blau sein wird.
Daher beschließt der neue Softwareentwickler herauszufinden, welche anderen Mitarbeiter auch Teil seines Teams sind.
Er macht also die Sucheingabe \textbf{Team Blau}.\\

Weiterhin möchte er genauer verstehen, wie die Entwicklung in dem Projekt abläuft.
Er beschließt sich herauszufinden, für welche Teile die unterschiedlichen Teams und Mitarbeiter zuständig sind.
Dazu macht der neue Softwareentwickler die Sucheingabe \textbf{Zuständigkeiten im Projekt}.

\section{Aufbau der Studie}
Im Rahmen der Bachelorarbeit soll eine Studie durchgeführt werden, welche die neue Suchfunktion mit der bestehenden Confluence-Suche vergleicht.
Bei der Darstellung der Ergebnisse der Studie dürfen keine Geschäftsgeheimnisse preisgegeben werden.
Die Ergebnisse der Studie dürfen also nur Zahlen enthalten, und nicht tatsächlich gefundene Informationen in der Wissensdatenbank.\\

Die Teilnehmer nehmen freiwillig an der Studie teil und entstammen aus genau einem Projekt der adesso SE.
Der Name des Projektes soll nicht benannt werden.
Für die Studie werden lediglich Softwareentwickler ausgewählt.
Bei der Einteilung der Softwareentwickler soll darauf geachtet werden, dass die Zeit, welche die Teilnehmer im Projekt verbracht haben, je Gruppe ähnlich ist.
Die Teilnahme an der Studie ist anonym.\\

Die Teilnehmer der Studie werden in zwei Gruppen eingeteilt.
Jeder Teilnehmer bekommt die Aufgabe eine spezifische Information herauszufinden.
Dabei wird die Zeit gemessen, wie lange der Teilnehmer gebraucht hat, um die Information zu finden.
Außerdem wird dokumentiert, ob die Information korrekt ist.
Zusätzlich wird die Anzahl der gemachten Sucheingaben dokumentiert.
Es werden insgesamt folgende Daten gesammelt und ausgewertet:
-	TeilnehmerID
-	AufgabenID
-	Suchfunktion
-	Benötigte Zeit
-	Anzahl von Sucheingaben
-	Korrektheit der Antwort\\

Das Problem bei einer AB-Studie könnte sein, dass die Teilnehmer unterschiedliche Kenntnisse über die Wissensdatenbank haben, sodass einige Teilnehmer besser darin sind Informationen zu finden als andere.
Befinden sich in einer Gruppe eher Teilnehmer, welche bessere Kenntnisse über die Wissensdatenbank haben, und in der anderen Gruppe Teilnehmer, welche eher schlechtere Kenntnisse haben, dann werden die Ergebnisse der Studie verfälscht.
Mithilfe einer AB/BA-Studie könnte dieses Problem gelöst werden.
Dabei muss darauf geachtet werden, dass die Teilnehmer bei der Verwendung der zweiten Suchfunktion nicht die gleichen Informationen finden müssen, wie bei der ersten Suchfunktion.
Denn nachdem sie die Informationen mit der ersten Suchfunktion gefunden haben, wissen sie bereits, wo sich die Information befindet.
Es muss also zwei verschiedene Datensets geben.
Um Ungleichheiten zwischen den Datensets zu vermeiden, soll das Ursprungsdatenset randomisiert in zwei gleich große Subdatensets aufgeteilt werden.
Die erste Gruppe bekommt das erste Datenset für die Confluence-Suche und das zweite Datenset für die neue Suchfunktion.
Die zweite Gruppe bekommt das erste Datenset für die neue Suchfunktion und das zweite Datenset für die Confluence-Suche.\\

Weiterhin kann es bei den Teilnehmern zu Lerneffekten kommen, sodass die Verwendung der jeweils zweiten Suchfunktion besser funktioniert als die Verwendung der jeweils ersten Suchfunktion.
Der Lerneffekt der Teilnehmer soll gemessen werden und bei der Auswertung beachtet werden.


\section{Auswertung der Ergebnisse}
TODO: Ergänzen


\section{Diskussion des Studienaufbaus}
TODO: Ergänzen

\chapter{Zusammenfassung und Ausblick}

\section{Zusammenfassung}
TODO

\section{Ausblick}
TODO




% ########################### BIBLIOGRAPHY ############################

{
\sloppy
\printbibliography[title={\label{chap:references}Literaturverzeichnis}]{} 
}

% ############################# APPENDIX ##############################

\part*{Anhang}
\newcounter{chapterAppendix}
\setcounter{chapterAppendix}{1}
\renewcommand{\thechapter}{\Alph{chapterAppendix}}
\chapter{Anhang}
\label{chap:appendix}

\section{TODO}
TODO

% ############################ DECLARATION ############################

% declaration of academic honesty / Eidesstattliche Erklärung
\chapter*{Eidesstattliche Erklärung}
\label{chap:declaration}

Hiermit versichere ich, dass ich die vorliegende Arbeit ohne Hilfe Dritter und nur mit den angegebenen Quellen und Hilfsmitteln angefertigt habe. Ich habe alle Stellen, die ich aus den Quellen wörtlich oder inhaltlich entnommen habe, als solche kenntlich gemacht. Diese Arbeit hat in gleicher oder ähnlicher Form noch keiner Prüfungsbehörde vorgelegen.
\vspace*{1cm}
\begin{flushright}\small \underline{\documentLocation, den \documentDate \hspace{0.4\textwidth}}\end{flushright}

\clearpage

% #####################################################################

\end{document}